\section{Data and Measurement}
\label{sec:data}

\subsection{What GitHub records}

GitHub hosts software projects and provides tools for developers to coordinate work. A \emph{repository} is a project's collection of files and their version history. Git, the version-control system used by GitHub, allows a developer to save a \emph{commit}---a recorded set of file changes---on a local computer without an Internet connection. A \emph{push} uploads one or more commits to a remote repository and requires connectivity. A push can therefore occur later than the work represented by its commits. Neither a commit nor a push measures hours worked or the economic value of the code.

A \emph{pull request} proposes changes for review; closing it need not mean acceptance, whereas merging it incorporates the proposed changes. An \emph{issue} records a task, question or bug. Reviews and comments record coordination. A \emph{release} publishes a named software version. Repository, branch or tag creation is also recorded as an event. GH Archive stores public GitHub event records and their publication times; it is not a record of all computer work, private development, or every local Git operation \citep{gharchive,githubevents}.

\subsection{GitHub users and country assignment}
\label{sec:users}

The analysis begins from GitHub accounts returned by country-location searches for Bangladesh and fourteen potential donor countries: the Philippines, Pakistan, Sri Lanka, Nepal, Indonesia, Vietnam, Malaysia, Thailand, Kenya, Nigeria, Ghana, Morocco, Jordan, and Tunisia. GitHub profile location is self-reported. I therefore interpret the sample as accounts associated with each country rather than verified physical residence at every point in time.

The country universe is the set for which the original project assembled location-search lists and a common historical event panel. It is a data-availability sample spanning South and Southeast Asia, Africa and the Middle East, not a random sample of countries or an automatically selected set of Bangladesh's closest economic peers. The archived materials do not establish an ex ante income, climate or electricity threshold that determined these fourteen countries, so I do not retrospectively claim one. All fourteen enter the optimization before weights are estimated; only Pakistan, Ghana and Kenya receive positive primary weights. \Cref{tab:advisor_donors} reports every country's sample size, pre-period activity and fitted weight. Small samples, country-specific disturbances and cross-border spillovers are evaluated as limitations rather than screened using the size of the treatment-period gap.

For each account, I use GH Archive to recover public daily GitHub events from May 1 through September 4, 2024; all event dates are in UTC. The activity requirement is that an account records any of the nine event types retained by the extraction query, including issue comments and commit comments not counted in \texttt{productive\_events}, on at least four days between May 1 and July 17. This requirement is necessary but not sufficient for inclusion: the further exclusions described below also apply. The final eligible sample contains 59,158 users, including 10,241 accounts associated with Bangladesh. This restriction focuses the design on developers with observable pre-shutdown activity and reduces the influence of dormant accounts. It also means that the estimand is an effect on pre-active public GitHub users, not on every person who has a GitHub account.

The sample also excludes suspected automation accounts, but the original login-pattern list and behavioral thresholds used for that step are unavailable. Applying the activity requirement alone yields 60,729 candidates, compared with 59,158 retained accounts; the 1,571 additional exclusions cannot be attributed to specific rules, because the original exclusion rules remain unresolved. Both stored bot flags are false for every retained account, which confirms their status but does not recover the exclusion rule. The stored behavioral metrics---maximum daily commits between May 1 and July 17, the number of dates with positive commits, and that count divided by 78 calendar days---can be reproduced for all retained accounts, but their exclusion thresholds remain unknown.

\subsection{Outcomes}
\label{sec:outcomes}

The primary outcome is the daily active-user rate: the share of eligible users in a country who record at least one productive GitHub action on day $t$. Operationally the indicator is $1\{\texttt{productive\_events}>0\}$. I prefer this extensive-margin measure for two reasons. First, GitHub activity is sparse at the user-day level; \citet{kreitmeir2024} similarly convert daily GitHub actions into binary indicators because such actions are relatively rare at the user level. Second, raw commit counts in this sample contain extreme outliers. The original extraction query, recovered from the project's query history, defines \texttt{productive\_events} as the count of seven event types: pushes, pull-request events, pull-request reviews, review comments, issue events, releases, and creation of a repository, branch or tag. Pull-request and issue events count across actions, not only openings or closures. Ordinary issue comments and commit comments do not enter this primary count, although they enter the collaboration measure. An independent re-extraction exactly reproduces the primary counts on all 799,417 overlapping eligible user-day records through July 23. The earlier component list omitted creation and several actions, explaining why a positive primary indicator need not coincide with those narrower components. The primary label describes this operational basket and should not be read as a judgement that every event produces useful software.

Let $\mathcal U_c$ be the fixed set of eligible accounts in country $c$, $N_c=|\mathcal U_c|$, and $P_{it}$ the count of these seven event types published by user $i$ on UTC date $t$. The outcome is
\begin{equation}
 Y_{ct}=100\,\frac{\sum_{i\in\mathcal U_c}\mathbf1\{P_{it}>0\}}{N_c}.
 \label{eq:primary_outcome}
\end{equation}
An eligible user with no event on a date contributes zero; the denominator is not the number of users observed that day. Several qualifying events still count as one active user-day. A change from 15 to 2 percent is a decline of 13 percentage points, while the relative change is $100(2-15)/15$ percent. Event publication dates and Git author timestamps are different measurements throughout the paper.

I construct secondary outcomes along the software-production pipeline. The secondary productive-output indicator indicates a day with commits, issue closure, pull-request closure, or a release. Shipping activity indicates a day with a merged pull request or a release. Collaboration activity indicates a day with code review or comment activity. I also construct count outcomes per 1,000 eligible users: collaboration events and shipping output. These measures separate recorded output from downstream delivery and coordination. The primary outcome and the secondary productive-output indicator are distinct measures: their Bangladesh pre-period means are 17.25 and 15.38 percent, respectively (\cref{tab:data}).

A small number of accounts generate implausibly large daily or weekly commit counts, so the primary country-day synthetic-control analysis does not rely on a raw commit-count headline. National holidays are coded at the country-day level and used in robustness checks.

\begin{table}[htbp]
\centering
\small
\caption{Data and outcome construction}
\label{tab:data}
\begin{threeparttable}
\begin{tabular}{@{}p{0.22\textwidth}p{0.44\textwidth}p{0.26\textwidth}@{}}
\toprule
Component & Definition / sample & Scale \\
\midrule
Primary outcome & At least one of the seven public event types defined in the text & Country-day percentage (effects in pp) \\
Secondary productive output & Any commits, issues closed, pull requests closed, or releases & User-day indicator / country-day rate \\
Shipping activity & Any merged pull request or release & User-day indicator / country-day rate \\
Collaboration activity & Any PR review, issue comment, PR review comment, or commit comment & User-day indicator / country-day rate \\
Collaboration events & Total review/comment events & Per 1,000 eligible users \\
Shipping output & Merged pull requests plus releases & Per 1,000 eligible users \\
Mechanism sample & Public repo-SHA pairs from fully observed PushEvents, matched to commit metadata & Commit level \\
\bottomrule
\end{tabular}
\begin{tablenotes}[flushleft]
\footnotesize
\item \textit{Notes:} Public GitHub events do not measure private repository activity, work on other platforms, meetings, or local code that is never pushed publicly.
\end{tablenotes}
\end{threeparttable}
\end{table}

\subsection{Commit-level mechanism data}
\label{sec:mechanism_data}

To distinguish coding from publication, I build a separate commit-level dataset for Bangladesh. I extract public PushEvents by eligible Bangladeshi users, retain PushEvents for which the historical payload reveals the complete set of commit SHAs, and keep the earliest observed push for each repository-SHA pair. I then query GitHub commit metadata to recover \texttt{authoredDate}, \texttt{committedDate}, and GitHub-linked author and committer logins. The metadata-query target contains 86,065 repository-SHA pairs, rather than the full pre-period universe: only 17,000 of 134,479 pairs first pushed before July 18 (12.64 percent) were queried, while all 69,065 pairs first pushed on or after July 18 were targeted; 70,576 commit objects are successfully resolved, an 82.0 percent resolution rate.

The strict mechanism sample requires (i) a non-missing author timestamp, (ii) a nonnegative author-to-push delay, and (iii) equality between the GitHub login linked to the commit author and the PushEvent actor. This leaves 59,920 commits. The equality condition avoids attributing a commit written by another person to the Bangladeshi user who happened to push it. Git author and committer timestamps are embedded metadata rather than independently verified work-time measurements, so they are interpreted as evidence consistent with offline production rather than as an infallible clock of when coding occurred.

\subsection{International repository-domain cohorts}
\label{sec:international_sector_data}

To construct international comparisons within repository domains, I retrieve the repository associated with each of the nine retained public event types for the unchanged 59,158 eligible accounts during May 1--July 23. For each account, I select its repository with the largest number of pre-shutdown days containing a primary event. Ties are resolved by the number of those events and then the repository name in lexical order. The rule uses May 1--July 17 only and selects a repository for 59,127 accounts; the remaining 31 have no qualifying primary event in that period. The selected repositories comprise 58,122 distinct projects. This is a mutually exclusive assignment, unlike the multi-sector activity measures used in the older Bangladesh recovery analysis.

I apply the same published multilingual NAICS classifier across all countries, using repository name and a root README from the latest retrievable default-branch commit no later than July 17, 2024. The model revision is pinned and the training project's text-formatting function is used. Historical source text is retrieved through GitHub using the same cutoff and procedure for Bangladesh and every donor country. Of the 58,122 selected repositories, 35,160 yield a nonempty historical README; 420 have empty README files, 11,804 have no root README, 9,949 are unavailable, 788 have no retrievable commit at the cutoff, and one has an unreadable text object. Current descriptions, topics and profile text are excluded. The main international assignment requires both an available historical README and a classifier score of at least 0.80. Accounts that fail these requirements remain in an explicit unclassified group. Sensitivities use 0.70 and 0.90 and, separately, permit name-only classifications at 0.80. Coverage is reported for every country in \cref{tab:sector_coverage}.

These groups measure a user's pre-treatment association with a repository application domain, not verified industry of employment. The most active repository need not represent all of a user's work. Retrieval follows the default branch available at query time, restricting its commit history to the pre-treatment cutoff; it cannot establish which branch was the default in July 2024. Histories may omit projects that became unavailable or whose history was rewritten, and the model's score is not a validated probability of correctness in this country sample. The reported test performance in the model card uses its own labelled data, not a hand-audited sample of these historical repositories \citep{quispe2026naics}. The design therefore addresses heterogeneity in observable activity across inferred domains, conditional on classification and coverage, rather than estimating industrial value added.

\subsection{Sector-recovery sample and repository classification}
\label{sec:sector_data}

To study the period after connectivity returned, I construct a separate within-Bangladesh panel. It starts from the 23,866 Bangladesh-associated accounts with at least one day of selected public GitHub activity between January 18 and July 17, 2024, and draws a fixed sample of 20,000 of them. Eligibility depends only on pre-shutdown activity, so the sample composition cannot respond to the shutdown or to post-shutdown behavior. For these accounts, I recover all public GH Archive events from January 18 through November 23, 2024, aggregated to the user-repository-day level.
This sample differs from the 59,158-account synthetic-control sample in its eligibility rule (one rather than four pre-treatment active days), its observation window, and its restriction to Bangladesh, so its activity levels are not directly comparable with the country-day rates in \cref{sec:results}.

Each repository is assigned to a two-digit sector of the North American Industry Classification System (NAICS) by a publicly released multilingual transformer classifier that predicts a NAICS sector from repository text (\texttt{aquiro1994/naics-github-classifier-multilingual}; \citealp{quispe2026naics}). The classifier input is the repository name and a historical snapshot of the root README file. For repositories observed on or before July 17, 2024, the snapshot is taken as of July 17; for repositories first observed after the shutdown, it is taken as of their first observed activity date. Current repository descriptions and topics are deliberately excluded, because they may have been written after the study period. Of the 168,875 distinct repositories in the panel, a historical README is recovered for 92,979 (55.1 percent); the remainder lack a root README at the target date or could not be retrieved. The main specification assigns a sector only when the classifier's confidence score is at least 0.80, which holds for 32,491 repositories across all sectors; repositories below this threshold do not contribute to any sector. Thresholds of 0.70 and 0.90 are used as robustness checks.

The sector measures the economic sector associated with the repositories in which observable GitHub activity occurs. It is not the developer's employer or industry of employment: the same developer may contribute to repositories classified in several sectors, and a repository's sector describes the application domain inferred from its text rather than the organization that funds the work.
